\chapter{Introduction}
\label{ch:introduction}

\setlength{\parskip}{0.7\baselineskip}

\section{Introduction}

In today's increasingly digital world, the internet has become an indispensable resource for communication, commerce, and information. However, this widespread connectivity has also given rise to a growing threat landscape of scam websites, phishing pages, and fraudulent online platforms that deceive and exploit unsuspecting users. Traditional approaches to identifying such threats — such as manually checking URLs or relying on browser warnings — are often insufficient, inconsistent, and inaccessible to non-technical users.

To address these limitations, we propose \textbf{ProblySUS}, a full-stack, multi-signal scam and phishing detection system that analyzes website URLs and produces a comprehensive, explainable risk assessment. The system integrates multiple techniques including domain intelligence, DNS analysis, threat feed lookups, content inspection, headless browser simulation, privacy auditing, and tracker identification to classify websites as \textbf{Safe, Caution, Suspicious, or Fraudulent} through a composite risk score ranging from 0 to 100.

By combining these signals into a single, accessible platform — available both as a web application and a browser extension — ProblySUS serves as an intelligent, real-time security companion for everyday internet users.

\subsection{Methodology}

The development of ProblySUS followed a structured and iterative methodology, beginning with a thorough requirement analysis. To understand the nature of the problem, we studied common phishing and scam patterns, reviewed existing detection tools, and examined open threat intelligence resources. This revealed that no single signal is sufficient for reliable scam detection — effective identification requires the combination of multiple independent indicators across URL structure, domain metadata, network behavior, and page content.

Based on these insights, we proceeded to technology stack selection. The backend was developed using the \textbf{Python / Flask} framework, chosen for its simplicity, flexibility, and rich ecosystem of analysis libraries. For domain analysis, we integrated \texttt{python-whois} and \texttt{dnspython}; for content inspection, \texttt{beautifulsoup4} and \texttt{textstat}; and for real browser simulation, \textbf{Playwright with headless Chromium}, enabling the interception of JavaScript-driven redirects and runtime network requests that static HTTP analysis cannot detect.

The frontend was built using \textbf{React 19 with Vite 7} and styled with \textbf{Tailwind CSS v4}, with \textbf{Framer Motion} for animations and \textbf{Axios} for API communication. Additionally, a \textbf{Chrome browser extension (Manifest V3)} was developed to bring one-click scanning directly into the user's browsing context.

In the system architecture design phase, we adopted a \textbf{modular pipeline approach}, wherein each analysis concern — URL validation, domain age, DNS records, URL patterns, content trust, blacklist lookup, browser behavior, network profiling, privacy grading, and tracker identification — is handled by a dedicated independent module. All module outputs are aggregated and passed to a central risk scoring engine (\texttt{scorer.py}) that applies a confidence-weighted penalty system with compound signal amplification and hard overrides for critical threats.

The architecture also incorporates a \textbf{ThreadPoolExecutor-based parallel execution model} that significantly reduces analysis time by running independent checks concurrently, and an \textbf{in-memory result cache with a 30-minute TTL} to avoid redundant scans.

The development and integration process was iterative. The backend API, frontend panels, and browser extension were developed and tested in parallel, with RESTful JSON serving as the integration contract. Once development was complete, we entered the testing and evaluation phase, which included unit testing of individual modules, integration testing of the full analysis pipeline, and performance testing to validate responsiveness under different network conditions.

A \textbf{Fast Mode (\texttt{?fast=true})} was also introduced to allow analysis in \textbf{3–5 seconds} by skipping the Playwright browser simulation, compared to the full \textbf{8–12 second deep scan}.

\subsection{Features}

ProblySUS is a multi-layered, cross-platform URL security analysis system designed to make scam detection accessible, transparent, and reliable for all users. The system performs over ten independent checks on any submitted URL and presents results through an intuitive, visually rich dashboard.

The system's threat intelligence backbone cross-references submitted domains against a local copy of the \textbf{URLhaus threat feed} — a community-maintained database of malicious URLs — that is automatically refreshed every 24 hours without manual intervention. A blacklist hit triggers an immediate \textbf{Fraudulent verdict}, bypassing further analysis.

For domains not found on the blacklist, the system evaluates \textbf{domain age via WHOIS lookup}, flags newly registered domains (under 30 days) as high-risk phishing signals, and inspects \textbf{DNS records} for the presence of MX and SPF entries that legitimate businesses are expected to maintain.

URL pattern analysis detects \textbf{IP-based hostnames, abusive top-level domains (.tk, .xyz, .pw), excessive hyphens, and embedded phishing keywords} — all common hallmarks of fraudulent sites. Content analysis fetches the live page and scans for urgency language, sensitive data requests, and the absence of trust indicators such as \textbf{Privacy Policy} and \textbf{Contact pages}.

Going beyond static analysis, ProblySUS launches a \textbf{headless Chromium browser via Playwright} to simulate a real user visit, intercepting all HTTP responses to detect redirect chains and all outgoing requests to profile the site's external network footprint. A dedicated \textbf{network intelligence module} classifies each contacted domain as safe infrastructure (CDN/cloud), suspicious, or unknown.

A \textbf{privacy auditing module} measures cookie tracking intensity, third-party script loading, and browser fingerprinting techniques (Canvas, WebGL, AudioContext, WebRTC, etc.), assigning a privacy grade from \textbf{Good} to \textbf{Invasive}. Finally, a \textbf{tracker identification module} matches loaded resources against a local database to identify named commercial trackers such as \textbf{Google Analytics, Facebook Pixel, and Hotjar}.

All findings are surfaced through a detailed \textbf{results dashboard} with animated panels, a risk score gauge, color-coded labels, and an explainable reasons list — as well as through the browser extension's compact popup, which presents the same data in a tabbed interface directly within the browser.

\subsection{Significance of Work}

ProblySUS addresses a critical and growing need in the digital age: empowering everyday users to make informed trust decisions about the websites they visit. With online scams, phishing attacks, and fraudulent e-commerce sites becoming increasingly sophisticated, traditional single-signal detection methods — such as SSL badge checks or basic blacklist lookups — are no longer adequate.

This project delivers a significant advancement by combining ten independent analysis signals into a unified, explainable risk score, reducing both false positives and false negatives through compound signal weighting.

The use of a \textbf{headless browser simulation} is particularly significant, as it enables the detection of JavaScript-driven redirects and dynamically loaded malicious resources that no static HTTP-based scanner can observe. The integration of \textbf{privacy and tracker auditing} further extends the system's utility beyond pure scam detection, providing users with a holistic view of a site's data collection practices.

The system's \textbf{local-first architecture} — with no third-party API keys required and a locally cached threat feed — ensures consistent performance, full privacy of submitted URLs, and zero dependency on paid external services.

The key contributions of ProblySUS include:

\begin{itemize}
\item \textbf{Multi-signal composite scoring} — 10+ independent analysis modules aggregated into a single 0–100 risk score with explainable reasons.
\item \textbf{Runtime behavior analysis} — Playwright headless browser captures JavaScript redirects and asynchronous network requests invisible to static analyzers.
\item \textbf{Privacy and tracker auditing} — measures fingerprinting intensity and identifies named commercial trackers, providing transparency beyond basic scam detection.
\item \textbf{Browser extension} — brings one-click URL scanning into the active browsing session with a tabbed detail view and deep-link to the web dashboard.
\end{itemize}